% Generated draft: requires mathematical and source review.
\documentclass[11pt]{article}
\usepackage[margin=25mm]{geometry}
\usepackage{fontspec}
\setmainfont{Latin Modern Roman}
\usepackage{amsmath,amssymb,mathtools,unicode-math}
\setmathfont{Latin Modern Math}
\usepackage{enumitem,needspace,xurl,hyperref}
\setlength{\parindent}{0pt}
\setlength{\parskip}{5pt}
\setlength{\emergencystretch}{3em}
\newcommand{\N}{\mathbb{N}}
\newcommand{\Z}{\mathbb{Z}}
\newcommand{\Q}{\mathbb{Q}}
\newcommand{\R}{\mathbb{R}}
\newcommand{\C}{\mathbb{C}}
\newcommand{\F}{\mathbb{F}}
\newcommand{\A}{\mathbb{A}}
\newcommand{\PP}{\mathbb{P}}
\newcommand{\E}{\mathbb{E}}
\newcommand{\eps}{\varepsilon}
\newcommand{\dd}{\,\mathrm{d}}
\begin{document}
\section*{UnsolvedMath ID: 30004287}
\textbf{Source problem code:} catalogue entry
\subsection{Definitions and mathematical statement}
The supplied statement does not specify the ground field. The formulation below explicitly adopts the usual number-field interpretation of the Hasse principle and Brauer--Manin obstruction; this domain is an interpretation requiring confirmation, not a hypothesis explicitly present in the source.

Let \(k\) be a number field, let \(\bar{k}\) be an algebraic closure, and let \(\Omega_k\) be the set of all places of \(k\), including the archimedean places. Write \(k_v\) for the completion at \(v\). Let \(X\) be a smooth, proper, geometrically integral variety over \(k\). Here geometrically integral means that \(X_{\bar{k}}\) is integral. Interpret rational connectedness geometrically: there is a dense open subset \(U\subseteq X_{\bar{k}}\) such that, for any \(x,y\in U(\bar{k})\), there is a morphism \(f:\PP^1_{\bar{k}}\to X_{\bar{k}}\) whose image contains both \(x\) and \(y\).

Let \(\A_k\) denote the adele ring of \(k\). Since \(X\) is proper, its adelic points can be identified with
\[
X(\A_k)=\prod_{v\in\Omega_k}X(k_v).
\]
Define the cohomological Brauer group by
\[
\operatorname{Br}(X)=H^2_{\mathrm{\acute{e}t}}(X,\mathbb G_m)_{\mathrm{tors}}.
\]
For \(\alpha\in\operatorname{Br}(X)\) and \(P_v\in X(k_v)\), evaluation gives \(\alpha(P_v)\in\operatorname{Br}(k_v)\). With \(\operatorname{inv}_v:\operatorname{Br}(k_v)\to\Q/\Z\) the local invariant map, define
\[
X(\A_k)^{\operatorname{Br}}=
\left\{(P_v)_v\in X(\A_k):
\sum_{v\in\Omega_k}\operatorname{inv}_v\bigl(\alpha(P_v)\bigr)=0
\text{ for every }\alpha\in\operatorname{Br}(X)\right\}.
\]
For each fixed \(\alpha\) and adelic point, only finitely many summands are nonzero. Global reciprocity puts the diagonal image of \(X(k)\) inside this set. A failure of the Hasse principle means that \(X(k_v)\ne\varnothing\) for every \(v\), but \(X(k)=\varnothing\); a Brauer--Manin obstruction to rational points occurs when \(X(\A_k)\ne\varnothing\) but \(X(\A_k)^{\operatorname{Br}}=\varnothing\).

Under the stated number-field and geometric rational-connectedness interpretations, the conjectural assertion is: for every number field \(k\) and every smooth, proper, geometrically integral, rationally connected \(k\)-variety \(X\),
\[
X(\A_k)^{\operatorname{Br}}\ne\varnothing
\quad\Longrightarrow\quad
X(k)\ne\varnothing.
\]
Thus the Brauer--Manin obstruction is asserted to be the only obstruction to the Hasse principle. No assertion about weak approximation is included. The supplied curation identifies this assertion with the Colliot--Th\'el\`ene conjecture and labels it open; these are supplied attribution and status information, not independently verified conclusions.
\subsection{Short English statement}
Under the explicitly provisional number-field interpretation, every smooth, proper, geometrically integral, geometrically rationally connected variety with an adelic point orthogonal to its entire Brauer group should have a rational point. The source omits the ground field, so the intended domain requires confirmation. The supplied record labels the conjecture open.
\subsection{Sources}
\begin{itemize}
\item Local UnsolvedMath registry, permanent ID \(30004287\), title ``Brauer--Manin Obstruction for Rationally Connected Varieties,'' exact source problem code \texttt{catalogue entry}, in the Oberwolfach Reports Open Problems collection. The record supplies proposed year 2019 and no named proposer. Its literature assessment identifies the assertion as the Colliot--Th\'el\`ene conjecture.
\item Supplied source citation: \emph{Analytic Number Theory (2020)}, Oberwolfach Reports, pp.\ 3141--3205, DOI \texttt{10.4171/owr/2019/50}, \url{https://doi.org/10.4171/owr/2019/50}. This is also the literature-source link labelled ``original OWR report.'' The supplied background describes these reports as documenting mathematical workshops and their open-problem discussions. No problem-specific page or passage is supplied.
\item Supplied literature source: Bright--Browning--Loughran, ``rational points and zero-cycles,'' \url{https://academic.oup.com/imrn/article/2021/6/4255/5292483}.
\item Supplied literature source: Browning--Loughran--Sofos, ``Hasse principle for random Fano hypersurfaces,'' \url{https://annals.math.princeton.edu/2023/197-3/p03}.
\item Literature assessment checked at \texttt{2026-08-22}: this date and assessment are supplied dated curation, not independently verified. The supplied assessment labels the problem open, says that the Colliot--Th\'el\`ene conjecture remains conjectural in the stated generality, and reports results for important special families together with continued conditional use of the general statement. The background explicitly calls this dated literature triage derived from the supplied OWR curation and requests independent verification of time-sensitive status and cited claims.
\end{itemize}
\end{document}
